% ID: SEL-D0027-OA-E03
% ESTADO: cotejada. Segunda pasada visual completa y borrador compilado revisados.
% Fuente: Selectividad/matematicas_aplicadas/suplemente_septiembre_EXAMEN 6-ANDALUCÍA.pdf
% SHA-256: 601fd1292e7d6b21ceecde2856a283e0e929d02d4e314d97fd7062bda846650a
% Página: 1; opción A; ejercicio original 3.

De los sucesos $A$ y $B$ se sabe que $P(A)=0.6$, $P(B/A)=0.8$ y $P(B/A^C)=0.1$.
\begin{enumerate}[label=\alph*)]
\item \textbf{(1.8 puntos)} Calcule las probabilidades $P(B)$, $P(A\cap B)$ y $P(A\cup B)$.
\item \textbf{(0.7 puntos)} \textquestiondown{}Son los sucesos $A$ y $B$ independientes?
\end{enumerate}
